\documentclass[10pt,a4paper]{amsart}
\usepackage[margin=19mm]{geometry}
\usepackage{amsmath,amssymb,mathtools}
\usepackage[scr=rsfs,cal=euler]{mathalfa}
\usepackage{xfrac}
\usepackage[unicode,hidelinks,bookmarks=false]{hyperref}
\newcommand{\ie}{i.e.\ }
\newcommand{\cf}{cf.\ }
\newcommand{\resp}{respectively\ }
\newcommand{\Subsection}{\subsection}
\providecommand{\nobreakdash}{\nobreak\hbox{-}}
\setlength{\emergencystretch}{2em}
\multlinegap0pt
\usepackage{bm}
\usepackage{bbm}
\usepackage{dsfont}
\usepackage{xspace}
\usepackage{cancel}
\usepackage{mathtools}
\usepackage{amsthm}
\makeatletter
\@ifundefined{theorem}{\newtheorem{theorem}{Theorem}}{}
\makeatother
\newcommand{\Mdash}{\acute{\mathcal{M}}}
\newcommand{\V}{\mathcal{V}}
\newcommand{\du}{d^\vee}
\newcommand{\im}{\textnormal{im}\,}
\newcommand{\inv}[0]{\textnormal{inv}}
\newcommand{\mf}[1]{\mathfrak{#1}}
\newcommand{\pl}[0]{\textnormal{pl}}
\newcommand{\quot}[0]{\textnormal{Quot}}
\newcommand{\res}{\textnormal{res}}
\newcommand{\rk}{\textnormal{rank}}
\newcommand{\vir}[0]{\textnormal{vir}}
\newcommand{\wh}[1]{\widehat{#1}}
\title[Joyce's invariant and Virasoro Constraints for Quot schemes ]{Joyce's invariant and Virasoro Constraints for Quot schemes on curves}
\author{Parvez Rasul}
\date{Source: arXiv:2608.04795v1 (CC BY 4.0)}
\begin{document}
\maketitle

\noindent\emph{Source and licence.} Joyce's invariant and Virasoro Constraints for Quot schemes on curves. Version arXiv:2608.04795v1 (CC BY 4.0), distributed under the Creative Commons Attribution 4.0 International licence, as recorded in the arXiv licence metadata for this exact version and shown on the abstract page https://arxiv.org/abs/2608.04795v1. Full rights notice: licenses/arxiv-2608.04795-CC-BY-4.0.txt.\par\smallskip

\noindent\textit{Editorial scope.} Counted unit: the Theorem whose source label is \texttt{theorem intersections on Quot N-1} in the article (source0.tex lines 2232--2244, source label \texttt{theorem intersections on Quot N-1}), delivered together with the article's own complete original proof of it (source0.tex lines 2245--2326, one \texttt{proof} environment of 82 lines). No mathematics is written, repaired or re-derived: statement and proof are the article's own text line for line. Required context retained uncounted: theorem invariant quot N-1 (lines 1352--1370); equation definition rho (lines 1362--1369). Every definition, display and label of those lines is retained unchanged. Omitted from this package and not redistributed: 1--1351, 1370--2231, 2327--2629. Nothing inside a retained excerpt is omitted or altered: every retained body is delivered exactly as the article's lines are written, so no omission receipt of any kind accompanies this package. Citations: the retained excerpts carry no citation command at all, so no bibliography entry of the article is transcribed and no reference is redistributed. Reference adaptation: none was needed and none was made; every reference inside the delivered unit resolves inside it, so no unresolved reference or citation is produced. Printed slips kept literal: no printed word, symbol, hypothesis or number of the counted statement or of its proof was altered. Layout and class: the article's own class is \texttt{amsart} with packages including mathabx, amsmath, amsfonts, amsthm, mathrsfs, amssymb, amscd, comment, enumerate, amsxtra, url, tikz-cd, physics, tcolorbox; the wrapper is the lane's pinned \texttt{amsart} editorial wrapper at 10pt on A4, which restates the theorem-like environments the retained text needs and adds the article's own macro definitions that the retained text needs (\texttt{\textbackslash Mdash}, \texttt{\textbackslash V}, \texttt{\textbackslash du}, \texttt{\textbackslash im}, \texttt{\textbackslash inv}, \texttt{\textbackslash mf}, \texttt{\textbackslash pl}, \texttt{\textbackslash quot}, \texttt{\textbackslash res}, \texttt{\textbackslash rk}, \texttt{\textbackslash vir}, \texttt{\textbackslash wh}), with extra wrapper packages (bm, bbm, dsfont, xspace, cancel, mathtools). The delivered numbering is the unit's own. Assets: no retained span contains a figure environment, a graphics-inclusion command, a PDF-page inclusion, a table or a TikZ picture; no image file is copied into this package, the package declares no asset dependency and no image file is redistributed with it. Licence: version arXiv:2608.04795v1 of this article is distributed under the Creative Commons Attribution 4.0 International licence, as recorded in the arXiv licence metadata for this exact version and shown on the abstract page https://arxiv.org/abs/2608.04795v1. A full rights notice is delivered at licenses/arxiv-2608.04795-CC-BY-4.0.txt. Characters: every retained excerpt is pure ASCII, so the delivered engine, XeLaTeX with its default Latin Modern text font, reports no missing character.
\begin{theorem}\label{theorem invariant quot N-1}
    Let $N =\rk(E)$ and $\nu = \chi(E)$.
    For any degree $d$, define $\du = d- \deg(E)$.
    The invariant 
    $$[Q_{\rk(E)-1,d}(E)]_\inv \quad  \in \quad \wh H_0(\Mdash_{(1,\du),1})/\im \acute T$$ is given by
    \begin{equation*}%\label{equation quot invariant}
    [Q_{\rk(E)-1,d}(E)]_\inv = 
    \res_{z=0} \left( \frac{1}{z^{\nu+N\du}} \cdot \rho(z) \cdot \sigma\left(\frac{N}{z} -s_{1,2,2}\right)\right) + \textnormal{im} \acute{T}\,,
    \end{equation*}
    where $\rho$ and $\sigma$ are defined as
\begin{align}\label{equation definition rho}
    \rho (z) & = \exp \biggl(
        \sum_{l=1}^\infty \frac{(-1)^l}{l!} \, z^l \, s_{{+},0,l}
    \biggr), \\
    \nonumber
    \sigma (x) & =
    \prod_{j=1}^g {} (x + s_{j,1,1} \, s_{j+g,1,1})\,.
\end{align}
\end{theorem}

\begin{theorem}\label{theorem intersections on Quot N-1}
    Let $\rk(E)=N$ and $\chi(E) = \nu$.
    Let $0 \leqslant s \leqslant g$ and $1 \leqslant j_1 < \cdots < j_s \leqslant g$ be integers
    such that the class $\prod_{i=1}^s b_1^{j_i}b_1^{j_i+g}$
    is nonzero.
    Then
    \begin{equation*}
        \int\limits_{[\quot_{N-1,d}(E))]^\vir} 
        \prod_{i=1}^s \left(b_1^{j_i}b_1^{j_i+g}\right) \cdot 
        a_1^{(\dim Q-s)} = N^{g-s}\,.
    \end{equation*}
    where $\dim Q = \nu + N(d - \deg(E)) - (1-g)$ is the expected dimension of $\quot_{N-1,d}(E)$.
\end{theorem}

\begin{proof}
We recall the setup.
Let $\du = d -\deg(E)$.
We have the translation operator
$$\acute T : H_{\bullet}(\Mdash_{(1,\du),1}) \to  H_{\bullet+2}(\Mdash_{(1,\du),1})\,.$$
Then Theorem \ref{theorem invariant quot N-1} says that the invariant 
    $[\quot_{N-1,d}(E))]_\inv$ is given by
    \begin{equation}\label{equation quot N-1 invariant}
    [\quot_{N-1,d}(E))]_\inv = 
    \res_{z=0} \left( \frac{1}{z^{\nu+N\du}} \cdot \rho(z) \cdot \sigma\left(\frac{N}{z} -s_{1,2,2}\right)\right) + \textnormal{im} \acute{T}\,,
    \end{equation}
    in $H_\bullet(\Mdash_{(1,\du),1})/\im \acute T$,
    where $\rho$ and $\sigma$ are defined as in \eqref{equation definition rho}.
Let 
$$\acute{T}^* : H^\bullet(\Mdash_{(1,\du),1}) \to H^{\bullet -2}(\Mdash_{(1,\du),1})$$ 
denote the dual of the map $\acute{T}$. 
The cohomology ring $H^\bullet(\Mdash^\pl_{(1,\du),1})$, 
being the dual of $H_\bullet(\Mdash^\pl_{(1,\du),1})$,
can be identified as $\ker (\acute{T}^*)$ inside $H^\bullet(\Mdash_{(1,\du),1})$.
So the induced pairing between $H_\bullet(\Mdash_{(1,\du),1})/\im \acute{T}$ and $\ker (\acute{T}^*)$
gives the pairing between
$H_\bullet(\Mdash^\pl_{(1,\du),1})$ and $H^\bullet(\Mdash^\pl_{(1,\du),1})$.
The map $\acute{T}^*$ is a derivation with the action on $S_{j,k,l}$ is given by
$$\acute{T}^* (S_{j,k,l}) = S_{j,k,l-1} \,.$$
%\qquad \text{ for }(j,k)\in \acute J \text{ and } l > k/2\,, $$
So in our case,  
\begin{equation}\label{elements of kernel Tdual}
(S_{1,2,2} - dS_{101}),\, (S_{j,1,1}S_{j+g,1,1}), \, 
(S_{+,0,1} - S_{1,0,1})
\in \ker(\acute{T}^*)\,.
\end{equation}
The classes $S_{j,k,l}$'s and $a_i, b_j^j, f_i$'s have the following relations,
\begin{equation}\label{relations of tautological classes}
    \mf h^* S_{1,0,1}=a_1, \quad  
    \mf h^* S_{j,1,1}=b_1^j, 
\quad \text{ and }  
\quad \mf h^* S_{1,2,2} = d^\vee a_1 - \sum_{j=1}^g {b_1^jb_1^{j+g}} \,.
\end{equation}
Also, note that the pullback of the universal sheaf $\V_1$ on the Quot scheme is trivial, which gives us 
$$ \mf h^*(S_{+,0,l}) = 0 \qquad \text{ for any } l>0 \,.$$
    Given $s$ and $j_1,\dots,j_s$ as in the statement of theorem,
    let us consider the class
    $$\zeta:= \prod_{i=1}^s \left(S_{j_i,1,1}S_{j_i+g,1,1} \right) \cdot
    (S_{1,0,1}-S_{+,0,1})^{\dim Q-s} \quad \in 
    H^\bullet(\Mdash_{(1,\du),1})\,.$$
    Using \eqref{elements of kernel Tdual}, we see that the class
    $\zeta$ is in $\ker(\acute{T}^*)$
    and using \eqref{relations of tautological classes},
    we have 
    $$\mf h^*(\zeta) = \prod_{i=1}^s \left(b_1^{j_i}b_1^{j_i+g}\right) \cdot 
        (a_1)^{\dim Q-s}\,.$$
    So the pairing of the cycle $[\quot_{N-1,d}(E))]_\inv$ and the class $\zeta$ gives us the required intersection, i.e.
    \begin{align*}
        \int\limits_{[\quot_{N-1,d}(E))]^\vir} 
        \prod_{i=1}^s  \left(b_1^{j_i}  b_1^{j_i+g}\right) \cdot (a_1)^{\dim Q-s} \
        = \left\langle  \zeta \,, \ [\quot_{N-1,d}(E))]_\inv \right\rangle\,.
    \end{align*} 
We simplify the residue appearing in the expression for 
$[\quot_{N-1,d}(E))]_\inv$ in
\eqref{equation quot N-1 invariant}. 
Let $\sigma_i$ be the coefficient of $(N/z)^i$ 
in the expression $\sigma(N/z-s_{1,2,2})$.
Clearly $\sigma_i$ is the
$i$-th symmetric polynomial in the variables
$$(s_{1,1,1}s_{1+g,1,1}-s_{1,2,2}), \cdots, (s_{g,1,1}s_{2g,1,1}-s_{1,2,2})\,.$$
Then the residue is equal to
\begin{align*}
    \sum_{k=0}^g \left( N^k \cdot \sigma_{g-k} \cdot \{z^{\nu+N(d-\deg(E))-1+k}\}(\rho(z)) \right)
\end{align*}
where $\{z^l\}(\rho(z))$ denotes the coefficient of $z^l$ in $\rho(z)$.
We write this expression as,
\begin{align*}
    \sum_{k=0}^g \left( N^k \cdot \sigma_{g-k} \cdot \{z^{\dim Q-g+k)}\}(\rho(z)) \right)
\end{align*}
So 
\begin{align*}
        \int\limits_{[\quot_{N-1,d}(E))]^\vir} 
        \prod_{i=1}^s & \left(b_1^{j_i}  b_1^{j_i+g}\right) \cdot (a_1)^{\dim Q-s}
         = \left\langle  \zeta \,, \ \sum_{k=0}^g \left( N^k \cdot \sigma_{g-k} \cdot \{z^{\dim Q-g+k)}\}(\rho(z)) \right) \right\rangle \,.
    \end{align*} 
This can be easily calculated to be $N^{g-s}$.
\end{proof}

\end{document}
